% ============================================================
%  Préambule — packages, mise en page, styles
% ============================================================

% --- Encodage & langue ---------------------------------------
\usepackage[utf8]{inputenc}
\usepackage[T1]{fontenc}
\usepackage[french]{babel}
\usepackage{lmodern}
\usepackage{textcomp}

% --- Mise en page --------------------------------------------
\usepackage[a4paper,left=2.5cm,right=2.5cm,top=2.5cm,bottom=2.5cm]{geometry}
\usepackage{setspace}
\onehalfspacing
\usepackage{parskip}
\setlength{\parindent}{1em}

% --- Graphiques & couleurs -----------------------------------
\usepackage{graphicx}
\graphicspath{{images/}}
\usepackage{xcolor}
\usepackage{tikz}
\usetikzlibrary{shapes.geometric,arrows.meta,positioning,calc,fit,backgrounds}

% --- Tableaux -------------------------------------------------
\usepackage{booktabs}
\usepackage{tabularx}
\usepackage{longtable}
\usepackage{multirow}
\usepackage{array}
\renewcommand{\arraystretch}{1.25}

% --- Maths ----------------------------------------------------
\usepackage{amsmath}
\usepackage{amssymb}

% --- Listes ---------------------------------------------------
\usepackage{enumitem}
\setlist{itemsep=2pt,parsep=0pt,topsep=4pt}

% --- Légendes -------------------------------------------------
\usepackage{caption}
\captionsetup{font=small,labelfont=bf,labelsep=endash}

% --- Palette --------------------------------------------------
\definecolor{accent}{HTML}{0E7C7B}
\definecolor{ink}{HTML}{1C2B39}
\definecolor{muted}{HTML}{5B6660}
\definecolor{soft}{HTML}{EEF1EF}
\definecolor{warn}{HTML}{C77D2E}
\definecolor{codebg}{HTML}{F7F8F7}
\definecolor{codekw}{HTML}{0B5FA5}
\definecolor{codestr}{HTML}{9C5A20}
\definecolor{codecom}{HTML}{6A7370}

% --- Code source ----------------------------------------------
\usepackage{listings}
\lstset{
  basicstyle=\ttfamily\footnotesize,
  keywordstyle=\color{codekw}\bfseries,
  stringstyle=\color{codestr},
  commentstyle=\color{codecom}\itshape,
  backgroundcolor=\color{codebg},
  frame=leftline,
  framesep=6pt,
  rulecolor=\color{accent},
  framerule=1.6pt,
  numbers=none,
  breaklines=true,
  breakatwhitespace=true,
  showstringspaces=false,
  columns=fullflexible,
  keepspaces=true,
  xleftmargin=10pt,
  aboveskip=10pt,
  belowskip=10pt,
  extendedchars=true,
  literate=%
    {é}{{\'e}}1 {è}{{\`e}}1 {ê}{{\^e}}1 {ë}{{\"e}}1
    {à}{{\`a}}1 {â}{{\^a}}1 {ä}{{\"a}}1
    {ù}{{\`u}}1 {û}{{\^u}}1 {ü}{{\"u}}1
    {î}{{\^i}}1 {ï}{{\"i}}1
    {ô}{{\^o}}1 {ö}{{\"o}}1
    {ç}{{\c c}}1
    {É}{{\'E}}1 {È}{{\`E}}1 {Ê}{{\^E}}1
    {À}{{\`A}}1 {Ç}{{\c C}}1
    {«}{{\guillemotleft}}1 {»}{{\guillemotright}}1
    {’}{{'}}1 {—}{{---}}1 {–}{{--}}1
}
\lstdefinestyle{python}{language=Python}
\lstdefinestyle{bash}{language=bash}
\lstdefinestyle{plain}{language=}

% --- En-têtes & pieds de page --------------------------------
\usepackage{fancyhdr}
\pagestyle{fancy}
\fancyhf{}
\fancyhead[L]{\small\nouppercase{\leftmark}}
\fancyhead[R]{\small\nouppercase{\rightmark}}
\fancyfoot[C]{\thepage}
\renewcommand{\headrulewidth}{0.4pt}
\renewcommand{\footrulewidth}{0pt}

\fancypagestyle{plain}{%
  \fancyhf{}%
  \fancyfoot[C]{\thepage}%
  \renewcommand{\headrulewidth}{0pt}%
}

\renewcommand{\chaptermark}[1]{\markboth{\chaptername~\thechapter.\ #1}{}}
\renewcommand{\sectionmark}[1]{\markright{\thesection.\ #1}}

% --- Titres de chapitre ---------------------------------------
\usepackage{titlesec}
\titleformat{\chapter}[display]
  {\normalfont\huge\bfseries\color{ink}}
  {\normalfont\normalsize\color{accent}\MakeUppercase{\chaptertitlename}\ \Huge\thechapter}
  {12pt}{\Huge}
\titlespacing*{\chapter}{0pt}{10pt}{28pt}

% --- Liens ----------------------------------------------------
% URL longues : autoriser la coupure aux separateurs pour qu'elles
% ne debordent pas de la marge dans la bibliographie.
\usepackage[hyphens]{url}
\Urlmuskip=0mu plus 1mu\relax
\def\UrlBreaks{\do\/\do\-\do\.\do\_\do\?\do\&\do\=\do\+\do\:}

\usepackage[hidelinks,breaklinks=true,plainpages=false]{hyperref}
\hypersetup{
  colorlinks=true,
  linkcolor=ink,
  citecolor=accent,
  urlcolor=accent,
  pdftitle={Rapport de stage - Systeme RAG},
  pdfauthor={\authorname}
}

% ============================================================
%  Commandes utilitaires
% ============================================================

% Valeur à compléter après exécution de l'évaluation RAGAS.
% Rend le trou visible à la relecture : remplacer par le chiffre mesuré.
\newcommand{\aremplir}[1]{\textcolor{warn}{\textbf{#1}}}

% Terme technique en anglais
\newcommand{\eng}[1]{\emph{#1}}

% Encadré "note"
\usepackage{mdframed}
\newmdenv[
  linecolor=accent,
  linewidth=1.6pt,
  topline=false,bottomline=false,rightline=false,
  backgroundcolor=soft,
  innerleftmargin=10pt,innerrightmargin=10pt,
  innertopmargin=8pt,innerbottommargin=8pt,
  skipabove=10pt,skipbelow=10pt
]{encadre}
